% answers_short.tex —— 2011 年·重庆卷·理科 · 简短答案速查
% 由 tools/make_short_answers.py 自 answers.tex 的 \gksolhead/\gkanswer 提取，
% 并与 2011/2011_数学真题_重庆理_答案要点.md 逐题比对后生成，勿手改（改源 MD 后重跑）。

\begin{gknotice}
\textbf{使用说明}：本节只给最终答案便于快速核对，完整推导见解析版。答案为编写者独立推导，\textbf{非官方标准答案}。
\end{gknotice}

\gkpart{选择题}
\begin{center}\begin{tabular}{ccccccccccc}
\toprule
题号 & 1 & 2 & 3 & 4 & 5 & 6 & 7 & 8 & 9 & 10 \\
\midrule
答案 & C & A & D & B & D & A & C & B & C & D \\
\bottomrule
\end{tabular}\end{center}

\gkpart{填空题}
\begin{itemize}[leftmargin=2.6em]
  \item[\textbf{11.}] $74$。
  \item[\textbf{12.}] $\sqrt3$。
  \item[\textbf{13.}] $\dfrac{11}{32}$。
  \item[\textbf{14.}] $-\dfrac{\sqrt{14}}2$。
  \item[\textbf{15.}] $\sqrt6-1$。
\end{itemize}

\gkpart{解答题}
\begin{itemize}[leftmargin=2.6em]
  \item[\textbf{16.}] 最大值 $2$，最小值 $\sqrt2$。
  \item[\textbf{17.}] (1) $\dfrac8{27}$；(2) $\xi$ 的分布列为 $P(\xi=1)=\dfrac1{27}$，$P(\xi=2)=\dfrac{14}{27}$，$P(\xi=3)=\dfrac49$；$E\xi=\dfrac{65}{27}$
  \item[\textbf{18.}] (1) $6x+2y-1=0$；(2) 极小值 $-3$（$x=0$），极大值 $15\mathrm e^{-3}$（$x=3$）
  \item[\textbf{19.}] (1) $\dfrac45$；(2) $\dfrac{\sqrt3}6$
  \item[\textbf{20.}] (1) $\dfrac{x^2}4+\dfrac{y^2}2=1$；(2) 存在，$F_1(-\sqrt{10},0)$，$F_2(\sqrt{10},0)$，$|PF_1|+|PF_2|=4\sqrt5$
  \item[\textbf{21.}] (1) $S_2=-2$，$a_3=\dfrac23$；(2) 证明见解析
\end{itemize}
